
\subsection{气体实验定律}
\begin{tabularx}{\columnwidth}{XX}
  玻意耳定律&$pV=C$\\
  查理定律（等压）&$p=CT$\\
  查理定律（等温）&$V=CT$\\
  理想气体状态方程&$pV=nRT$\\
  液体上下压强& $p_{\text{下}}-p_{\text{上}}=\rho g(h_{\text{下}}-h_{\text{上}})$\\
  有质量活塞上下压强& $p_{\text{下}}-p_{\text{上}}=\frac{Mg}{S}$\\
  水平液柱与活塞压强& $p_{\text{左}}=p_{\text{右}}$\\
  竖直液柱加速度& $p_{\text{下}}-p_{\text{上}}=\rho g_{等效}h$\\
  竖直活塞加速度& $p_{\text{下}}-p_{\text{上}}=\frac{M}{S}g_{等效}$\\
\end{tabularx}

\begin{formulanotebox}
  等效重力加速度，以a为加速度，加速上升时，$g_{等效}=g+a$，加速下降时，$g_{等效}=g-a$。
  液体压强规律：液体压强与液体密度、液体深度有关，与液体质量无关。
  气体压强规律：容器内气体压强处处相等。
  注意识别定律适用前提（一定质量气体）；温度计算需统一使用热力学温标（K）。
\end{formulanotebox}

\begin{keypointbox}
  \textbf{理想气体定义}：气体分子间没有相互作用力，分子体积可以忽略不计，只考虑分子热运动。
  常见的气体在高温低压下近似理想气体。\\
  \textbf{理想气体的内能与温度成正比。}\\
\textbf{过程识别优先}：先判断过程类型再选公式，避免在同一过程里混套不相容关系式。
\end{keypointbox}

\begin{casetypebox}
% TODO: 本节题型短标题未提供，请补充后逐题型填写
打气问题

气体罐装问题

液柱移动问题

\end{casetypebox}
